%% 6-15 まとめ：物理は「微分方程式＋はじめの状態」で書ける（第14回：組 V6-8。ページ制限なし。偶数ページ）
%% ほかの巻・物理の本との約束：
%%   5-5「運動の法則は『2回微分をふくむ式』。これが第6巻の微分方程式です（6-15）」、4-1「6-15 では、ばねも電気も熱も、少数の法則から同じ形の式で書けることを見ます」を回収。
%%   相似の表は物理第1巻の表2と同じ組（力↔電圧、速度↔電流、質量↔インダクタンス、ダンパー↔抵抗、ばねのやわらかさ 1/k↔静電容量）。
%%   1階の仲間は 6-4⑧の表、RL と落ちる物は 6-5⑦（質量↔L、b↔R、mg↔E、v↔i）、ばね・質量・ダンパー↔RLC は 6-6⑧・6-7⑦（m↔L、c↔R、k↔1/C、x↔q）。
%%   ★2〜3 の章なので、第3部の★5の内容は表とグラフで思い出すだけにし、解かせない。
\makeatletter
\newcommand\kviphead[1]{\par\addvspace{6pt}\Needspace{6\baselineskip}\noindent{\sffamily\bfseries\color{\bkTT}#1}\par\nobreak\vspace{2pt}}
\newcommand\kvipbox[1]{\par\addvspace{3pt}\noindent\textbf{#1}\par\vspace{3pt}\noindent\ignorespaces}
\newcommand\kvipline[1]{\par\addvspace{3pt}{\centering #1\par}\addvspace{3pt}\noindent\ignorespaces}
\newcommand\kvipframe[2][100mm]{\par\addvspace{4pt}{\centering\tikz\node[draw=\bkTT, line width=0.9pt, fill=bkTheme!10, rounded corners=1.5mm, inner xsep=6mm, inner ysep=2.5mm, align=center]{#2};\par}\addvspace{4pt}\noindent\ignorespaces}
\newcommand\kvipa{\par\vspace{2pt}\noindent\rightskip=0pt plus 8em\relax\hangindent=2.6\zw\hangafter=0\relax}
\newcommand\kvips{\rightskip=0pt plus 8em\relax}
% ばね（横）・台車・かべ（6-14 と同じ形）
\newcommand\kvipspring[2]{\draw[bkInk, line width=0.8pt, decorate, decoration={zigzag, amplitude=1.3mm, segment length=2.2mm, pre length=1.5mm, post length=1.5mm}] #1 -- #2;}
\newcommand\kvipwall[2]{\begin{scope}[shift={#1}]
  \draw[bkInk, line width=0.9pt] (0,0) -- (0,#2);
  \foreach \kvipy in {0,2.5,...,#2} {\draw[bkInk!60, line width=0.4pt] (0,\kvipy) -- (-1.6,{\kvipy+1.6});}
\end{scope}}
% 小さな3つの絵（ばね・回路・コップ）  \kvipspringpic{(x,y)}  \kviprlcpic{(x,y)}  \kvipcuppic{(x,y)}
\newcommand\kvipspringpic[1]{\begin{scope}[shift={#1}]
  \draw[bkInk, line width=0.8pt] (0,0) -- (40,0);
  \kvipwall{(1,0)}{14}
  \kvipspring{(1,10)}{(22,10)}
  \draw[bkInk, line width=0.8pt] (1,4) -- (9,4) (17,2) -- (9,2) -- (9,6) -- (17,6) (13,4) -- (22,4);
  \draw[bkInk, line width=1.4pt] (13,2.6) -- (13,5.4);
  \filldraw[fill=bkBlue!25, draw=bkBlue!80!black, line width=0.8pt, rounded corners=0.8mm] (22,0.6) rectangle (34,13);
  \node[font=\sffamily\small] at (28,6.8) {$m$};
  \node[font=\sffamily\footnotesize, anchor=south] at (11.5,12.6) {$k$};
  \node[font=\sffamily\footnotesize, anchor=north] at (13,1.4) {$c$};
\end{scope}}
\newcommand\kviprlcpic[1]{\begin{scope}[shift={#1}]
  \draw[bkInk, line width=0.8pt] (0,0) -- (0,4.5) (0,9.5) -- (0,14) -- (8,14);
  \draw[bkInk, line width=0.8pt] (0,7) circle (2.5);
  \node[font=\footnotesize] at (0,7) {$E$};
  \draw[bkInk, line width=0.8pt] (8,12) rectangle (17,16);
  \draw[bkInk, line width=0.8pt] (17,14) -- (20,14) \foreach \i in {0,...,2} { arc (180:0:1.5) } -- (30,14) -- (30,8.2) (30,5.8) -- (30,0) -- (0,0);
  \draw[bkInk, line width=1.3pt] (27,8.2) -- (33,8.2) (27,5.8) -- (33,5.8);
  \node[font=\sffamily\footnotesize, anchor=south] at (12.5,16.3) {$R$};
  \node[font=\sffamily\footnotesize, anchor=south] at (23.5,16.3) {$L$};
  \node[font=\sffamily\footnotesize, anchor=west] at (33.5,7) {$C$};
\end{scope}}
\newcommand\kvipcuppic[1]{\begin{scope}[shift={#1}]
  \fill[bkOrange!25] (2,0.8) -- (3.4,11) -- (16.6,11) -- (18,0.8) -- cycle;
  \draw[bkInk!80, line width=0.9pt] (1,14) -- (3,0) -- (17,0) -- (19,14);
  \draw[bkInk!80, line width=0.9pt] (18.4,10) .. controls (24,10) and (24,3) .. (17.6,3);
  \draw[bkGray, line width=0.7pt] (6,16) .. controls (5,18) and (7,19) .. (6,21) (10,16) .. controls (9,18) and (11,19) .. (10,21) (14,16) .. controls (13,18) and (15,19) .. (14,21);
\end{scope}}
\makeatother

\bkchapter{6}{15}{まとめ：物理は\\「微分方程式＋はじめの状態」で書ける}{\bklevelrange{2}{3}}{まとめ：物理は「微分方程式＋はじめの状態」で書ける}
\bkmichishirube{4-1,5-5,5-12,6-1,6-4,6-6,6-7,6-14}{7-0,8-13,9-5,9-10}

\begin{goal}
\begin{itemize}[label={\color{\bkTT}●}, leftmargin=1.6\zw, itemsep=0pt, topsep=1pt]
\item \bklevel{2}\enspace 運動の法則・回路のきまり・熱の出入りから、微分方程式とはじめの状態を書ける。
\item \bklevel{2}\enspace 相似の表で、ばね・電気・熱の量を対応させ、同じ形の答えを使える。
\item \bklevel{3}\enspace 式の形から、どの解き方（この巻のどの章）を使えばよいか選べる。
\end{itemize}
\end{goal}

\begin{hatena}
\noindent\begin{minipage}[c]{0.5\linewidth}
\raggedright
ばねにつけたおもり、電気の回路、冷めるお湯。この巻では、まったくちがう物なのに、同じ形の式が何度も出てきました。

なぜ、同じ形になるのでしょう。そして、同じ形だと、何がうれしいのでしょう。
\end{minipage}\hfill
\begin{minipage}[c]{0.46\linewidth}
\centering
\begin{tikzpicture}[x=0.8mm, y=0.8mm]
  \kvipspringpic{(0,27)}
  \kviprlcpic{(2,0)}
  \kvipcuppic{(42,28)}
  \node[draw=bkYellow!70!black, line width=0.8pt, rounded corners=1.5mm, fill=white, font=\sffamily\footnotesize, align=center, inner sep=2mm] (q) at (52,7) {同じ\\形の式？};
\end{tikzpicture}
\end{minipage}
\end{hatena}

\noindent この章は、第6巻のまとめです。新しい計算はほとんどありません。そのかわり、これまでの章をつなげて、「物理の法則は、なぜ微分方程式になるのか」「なぜ、ちがう物が同じ形の式になるのか」を考えます。最後に、この巻で学んだ解き方を1枚の地図にまとめます。

まとめの章なので、前の章を全部おぼえていなくても大丈夫です。わからない所は、章の番号（6-4 など）をたよりにもどって読めば、つながりが見えてきます。% Ver2 No.24 声かけ

%% ---------------- やってみよう・思い出そう ----------------
\begin{yattemiyou}[式のカードを仲間分けしよう]
この巻に出てきた式を、8枚のカードにしました。
\begin{bkcenter}
\begin{tblr}{width=\linewidth, colspec={Q[c,m,wd=6mm] X[l,m] Q[c,m,wd=6mm] X[l,m]}, rowsep=3pt, colsep=3pt, hlines={0.4pt,bkLine}, vlines={0.4pt,bkLine}, column{1,3}={font=\sffamily\bfseries}}
ア & 冷めるお湯 $\dfrac{\dd T}{\dd t}=-0.05(T-20)$（6-4） & イ & ばね $0.5\ddot x=-8x$（6-6） \\
ウ & RC の充電 $\dfrac{\dd V}{\dd t}=\dfrac{10-V}{2}$（6-4） & エ & LC 回路 $L\ddot q=-\dfrac{q}{C}$（6-6） \\
オ & 落ちる物 $m\dfrac{\dd v}{\dd t}=mg-bv$（6-5） & カ & ばね・ダンパー $m\ddot x+c\dot x+kx=0$（6-7） \\
キ & RL 回路 $L\dfrac{\dd i}{\dd t}=E-Ri$（6-4） & ク & RLC 回路 $L\ddot q+R\dot q+\dfrac{q}{C}=0$（6-7） \\
\end{tblr}
\end{bkcenter}
\begin{enumerate}[label={\bkcircled[bkGreen]{\arabic*}}, leftmargin=2.4\zw, itemsep=1pt, topsep=3pt]
\item 何回微分が入っているか（1回か2回か）で、2つのグループに分けよう。
\item 同じグループの中で、文字をかえるとぴったり同じ形になる組を見つけよう。
\item それぞれの組で、「何」と「何」が入れかわっているかを書き出そう。
\end{enumerate}
\end{yattemiyou}

\noindent 1回微分のグループは、ア・ウ・オ・キ。どれも「目標の値との差に比例して近づく」形（6-4④）です。2回微分のグループは、イ・エ・カ・ク。イとエは「ゆれる」形（6-6）、カとクは「ゆれながら止まる」形（6-7）です。

③で、たとえばオとキをくらべると、質量 $m$ とコイル $L$、抵抗の係数 $b$ と抵抗 $R$、重さ $mg$ と電池 $E$、速さ $v$ と電流 $i$ が入れかわっています（6-5⑦）。このような入れかわりの組を、この章で1つの表にまとめます。

\begin{omoidasou}
\bkrecall{4-1}{物理の\textbf{法則}は、実験と観察で見つけたきまり。数学の\textbf{定理}は、定義と公理から証明で導いたきまり。}
\bkrecall{5-5}{力の合計＝質量×加速度（$F=ma$）。加速度は位置の2回微分。}
\bkrecall{5-12}{ばねを $x$ のばす仕事は $\frac12kx^2$。コンデンサーにたまるエネルギーは $\frac12CV^2$。}
\bkrecall{6-1}{ルールの作り方：主役の量を決める→変わる速さを微分で書く→法則を式にして「＝」でつなぐ。はじめの状態は、階数の数だけ要る。}
\bkrecall{6-4}{目標へ近づく形 $\dfrac{\dd y}{\dd t}=\dfrac{y_\infty-y}{\tau}$。$\tau$ は時定数（RC、$L/R$、$C_{\mathrm{th}}R_{\mathrm{th}}$、$m/b$）。}
\bkrecall{6-7}{ばね・質量・ダンパー $m\ddot x+c\dot x+kx=0$。減衰比 $\zeta=\dfrac{c}{2\sqrt{mk}}$ で、ゆれ方が決まる。}
\end{omoidasou}

%% ---------------- 数学の道具 ----------------
\bksection{doguu}{数学の道具}

\bksubsection{1}{法則は「次の瞬間」を決める}

\label{bk:6-15:hosoku}\bkidx{ほうそく}{法則}5-5 で、「運動の法則は『2回微分をふくむ式』になります。これが第6巻の微分方程式です」と書きました。ニュートンの運動の法則
\kvipframe[80mm]{力の合計 $=m\times$ 加速度\qquad $F=m\dfrac{\dd^2x}{\dd t^2}$}
は、「今の位置と速度がわかれば、次の瞬間の速度の変わり方（加速度）がわかる」という形をしています。力がばねの力 $-kx$ なら $m\ddot x=-kx$（6-6）、重さと空気の抵抗なら $m\dot v=mg-bv$（6-5）。どちらも、力のきまりを入れただけで、微分方程式になりました。

\begin{bkfigure}
\begin{tikzpicture}[x=1mm, y=1mm]
  \node[bk node, minimum width=34mm, minimum height=13mm] (a) at (0,0) {今の状態\\（位置・速度）};
  \node[bk node, draw=bkOrange, minimum width=34mm, minimum height=13mm] (b) at (44,0) {法則\\（変わる速さを決める）};
  \node[bk node, minimum width=34mm, minimum height=13mm] (c) at (88,0) {次の瞬間の状態};
  \draw[bk link] (a) -- (b);
  \draw[bk link] (b) -- (c);
  \draw[bk link, bkBlue] (c.south) -- ++(0,-8) -| node[bk link label, pos=0.25, below] {くり返す（6-2 の表計算）} (a.south);
\end{tikzpicture}
\bkfigcaption{物理の法則は「今の状態から、次の瞬間の状態を決めるルール」。はじめの状態から、これをくり返すと未来が決まる。}
\end{bkfigure}

\kviphead{法則と、導いたきまり}
4-1 では、数学の「定理」（証明で導いたもの）と、物理の「法則」（実験で確かめたもの）を区別しました。この巻に出てきたきまりを、2つに分けてみます。
\begin{bkcenter}
\begin{bktable}{width=\linewidth, colspec={X[l,m] X[l,m]}, rowsep=2pt, row{1}={font=\sffamily\small}}
法則（実験で確かめられた出発点） & 導いたきまり（微分方程式を解いて出したもの） \\
運動の法則 $F=ma$ & ばねの周期 $T=2\pi\sqrt{m/k}$（6-6） \\
フックの法則 $F=-kx$ & 終端速度 $v_{\mathrm f}=mg/b$（6-5） \\
オームの法則 $V=RI$ & 時定数 $\tau=RC$、$L/R$（6-4） \\
コンデンサー $q=CV$・コイル $V=L\dfrac{\dd i}{\dd t}$ & 電気振動の周期 $2\pi\sqrt{LC}$（6-6） \\
ニュートンの冷却の法則（近似） & 冷める温度 $T=20+60e^{-t/20}$（6-4） \\
\end{bktable}
\bktabcaption{左の少数の法則から、右のたくさんのきまりが、数学で導かれる。}
\end{bkcenter}
高校の物理で「覚える式」として出てくる右の列のきまりは、どれも左の法則から、微分方程式を解いて出したものです。4-1 で「少数の法則から同じ形の式で書けることを見ます」と書いたのは、このことです。

\begin{column}{ラプラスの悪魔}
1814年、フランスの数学者ラプラスは、「ある瞬間の、すべての物の位置と速度を知り、それをすべて計算できる知性があれば、その知性にとって未来はすべて見通せるだろう」と書きました。この「知性」は、のちに「ラプラスの悪魔」とよばれます。この章の「微分方程式＋はじめの状態で未来が決まる」を、宇宙全体にまで広げた考えです。けれども、6-14 のコラムで見たように、ほんの少しのはじめのずれが大きく育つ（カオス）ので、じっさいには、はじめの状態をどれだけ正確にはかれるかで、予想できる先の長さが決まります。
\end{column}

\bksubsection{2}{ためる・失う・押す：式を立てる3つの部品}

\label{bk:6-15:buhin}ばねも電気も熱も、同じ形の式になるのはなぜでしょう。式を立てたときの「部品」を、くらべてみます。

\begin{bkfigure}
\begin{tikzpicture}[x=1mm, y=1mm]
  \kvipspringpic{(0,0)}
  \kviprlcpic{(48,0)}
  \kvipcuppic{(100,0)}
  \node[font=\sffamily\footnotesize, anchor=north, align=center] at (18,-2) {ばね・質量・ダンパー\\（6-7）};
  \node[font=\sffamily\footnotesize, anchor=north, align=center] at (63,-2) {RLC 回路\\（6-7）};
  \node[font=\sffamily\footnotesize, anchor=north, align=center] at (110,-2) {冷めるお湯\\（6-4）};
\end{tikzpicture}
\bkfigcaption{3つのしくみ。どれも「ためるもの」と「失うもの」があり、それぞれのきまりが比例の形。}
\end{bkfigure}

\begin{itemize}[leftmargin=1.2\zw, itemsep=3pt, topsep=2pt, label={\color{\bkTT}●}]
\item \textbf{位置・電気・熱をためる部品（ためるもの）}\par\noindent ばねはのびの形でエネルギーをため（$\frac12kx^2$）、コンデンサーは電気をため（$\frac12CV^2$）、お湯は熱をためます（熱容量 $C_{\mathrm{th}}$）。
\item \textbf{運動の勢いをためる部品（慣性）}\par\noindent おもりは速さの形でエネルギーをため（$\frac12mv^2$）、コイルは電流の形でためます（$\frac12Li^2$）。「流れ続けようとする」はたらきです。
\item \textbf{失う部品}\par\noindent ダンパー、抵抗、熱抵抗。速さや電流や熱の流れに比例して、エネルギーを熱にして逃がします。
\item \textbf{押す部品}\par\noindent 外からの力、電池の電圧、部屋との温度差。
\end{itemize}
（かっこの中の名前は、③の相似の表の行の名前です。）そして、部品のきまりが、どれも「比例」です（ばねの力は $kx$、抵抗の電圧は $Ri$、熱流は温度差÷$R_{\mathrm{th}}$）。比例のきまりを「合計」でつなぐので、どの分野でも、定数がかかった微分の項をたした線形の式（5-18）になるのです。

\kviphead{部品を「合計」でつなぐきまり}
部品のきまりだけでは、式はできません。部品どうしを「合計」でつなぐきまりが、分野ごとに1つずつあります。
\begin{bkcenter}
\begin{bktable}{width=\linewidth, colspec={Q[l,m,wd=16mm] X[l,m] X[l,m]}, rowsep=2pt, row{1}={font=\sffamily\small}}
 & つなぐきまり & 式 \\
力学 & 力の合計＝質量×加速度（運動の法則） & $m\dfrac{\dd v}{\dd t}=F-cv-kx$ \\
電気 & 電池の電圧＝部品の電圧の合計（回路を1周すると電圧の上がり下がりの合計は0。キルヒホッフの電圧の法則） & $E=L\dfrac{\dd i}{\dd t}+Ri+\dfrac{q}{C}$ \\
熱 & たまる熱＝入る熱－出る熱（エネルギーの保存） & $C_{\mathrm{th}}\dfrac{\dd T}{\dd t}=P-\dfrac{T-T_{\mathrm a}}{R_{\mathrm{th}}}$ \\
\end{bktable}
\bktabcaption{3つの分野の「つなぐきまり」。右の列が、そのまま微分方程式になる。}
\end{bkcenter}
\bkidx{きるひほっふのでんあつのほうそく}{キルヒホッフの電圧の法則}
電気の行の「1周すると0」は、2-15・6-6・6-7 で「電圧が2つに分かれる」「同じ2本の線につながっているので電圧が同じ」として使ってきたことを、1つのきまりにしたものです。熱の行は、6-4⑥の電気ポットの式そのものです（$T_{\mathrm a}$ は部屋の温度）。

\kviphead{\bkhatenaref の答え}
ばねも電気も熱も、①「ためるもの・慣性・失うもの・押すもの」という同じ役わりの部品でできていて、②部品のきまりがどれも比例で、③それを「合計」でつなぐからです。だから、同じ形の式になります。うれしいのは、1つの答えを作れば、文字をおきかえるだけで全部に使えることです（④で確かめます）。

\bksubsection{3}{相似の表}

\label{bk:6-15:soji}ちがう分野なのに、式が同じ形になることを\textbf{相似}（そうじ）といいます。\bkidx{そうじ}{相似（式の形）}力学の式と電気の式をならべます。
\kvipline{力学：\quad $m\dfrac{\dd v}{\dd t}+cv+kx=F$\quad （$v=\dfrac{\dd x}{\dd t}$）}
\kvipline{電気：\quad $L\dfrac{\dd i}{\dd t}+Ri+\dfrac{q}{C}=E$\quad （$i=\dfrac{\dd q}{\dd t}$）}
左から1つずつ見ると、$m\leftrightarrow L$、$c\leftrightarrow R$、$k\leftrightarrow1/C$、$x\leftrightarrow q$、$v\leftrightarrow i$、$F\leftrightarrow E$ と対応しています（6-6⑧・6-7⑦）。物理シリーズ『力学』の表2と同じ組です。熱と、6-14 の水そうも入れて、1つの表にします。
\begin{bkcenter}
\small
\begin{tblr}{width=\linewidth, colspec={Q[l,m,wd=21mm] X[c,m] X[c,m] X[c,m] X[c,m]}, row{1}={bg=bkTheme!12, font=\sffamily\bfseries\small}, column{1}={font=\sffamily\bfseries\small}, hlines={0.4pt,bkLine}, rowsep=2pt, colsep=3pt}
役わり & 力学 & 電気 & 熱 & 水（管とタンク） \\
流れるもの & 速度 $v$ & 電流 $i$ & 熱流 $P$ & 流量 $Q$ \\
押すもの & 力 $F$ & 電圧 $E$ & 温度差 $\Delta T$ & 圧力差（水位の差） \\
ためるもの & ばねのやわらかさ $1/k$ & 静電容量 $C$ & 熱容量 $C_{\mathrm{th}}$ & タンク \\
慣性 & 質量 $m$ & インダクタンス $L$ & （ない） & 管の中の水 \\
失うもの & ダンパー $c$ & 抵抗 $R$ & 熱抵抗 $R_{\mathrm{th}}$ & 管の流れにくさ \\
\end{tblr}
\bktabcaption{相似の表（物理シリーズ『力学』の表2と同じ組）。力 $\leftrightarrow$ 電圧、速度 $\leftrightarrow$ 電流。物理シリーズでは、電源の電圧を小文字の $e$ と書く。}
\end{bkcenter}
\Needspace{8\baselineskip}
\begin{itemize}[leftmargin=1.2\zw, itemsep=3pt, topsep=2pt, label={\color{\bkTT}●}]
\item \textbf{ばねは「やわらかさ $1/k$」がコンデンサーの $C$。}\par\noindent ばねの力は $kx$、コンデンサーの電圧は $\dfrac{q}{C}=\dfrac1C\times q$。$k$ と $\dfrac1C$ が同じ役わりなので、「ためるもの」の列は $k$ ではなく $\dfrac1k$ を書きます。やわらかいばね（$k$ が小さい）ほど、たくさんのびてためられる、という意味です。
\item \textbf{熱には「慣性」がない。}\par\noindent 熱の流れには、質量やコイルのような「流れ続けようとする」部品がありません。だから熱の式は、ためる部品と失う部品だけの1階の式になり、ゆれる答え（6-6・6-12 の $\sin$・$\cos$）が出てきません。お湯の温度が、部屋の温度を行きすぎて下がり、またもどる、ということはないのです（6-14 の2つの物の温度も、ゆれずに近づきました）。
\end{itemize}

\begin{seikaku}
相似は「式の形が同じ」ということで、物そのものが同じなのではありません。たとえば力学と電気では、力×速度も電圧×電流も1秒あたりのエネルギー（仕事率、W）になりますが、熱では、熱流がすでに1秒あたりのエネルギーなので、温度差×熱流は仕事率になりません。また、力と電流を組にする、もう1つの相似の決め方もありますが、このシリーズでは物理第1巻の表と同じ「力↔電圧」にそろえます。
\end{seikaku}

\kviphead{1回微分の仲間：時定数と最後の値}
6-4⑧の表に、6-5 と 6-14 の行を加えて、もう一度ならべます。どれも $\dfrac{\dd y}{\dd t}=\dfrac{y_\infty-y}{\tau}$ の形で、答えは $y=y_\infty+(y_0-y_\infty)e^{-t/\tau}$ です。
\begin{bkcenter}
\small
\begin{tblr}{width=\linewidth, colspec={X[l,m] Q[c,m] Q[c,m] Q[c,m] Q[c,m]}, row{1}={bg=bkTheme!12, font=\sffamily\small}, hline{1,2,Z}={0.6pt,bkGray}, hline{3-8}={0.3pt,bkGray}, rowsep=2.5pt, colsep=4pt}
しくみ & 変わる量 & 最後の値 & 時定数 & 例の数 \\
冷めるお湯（6-4） & 温度 $T$ & 部屋の温度 & $C_{\mathrm{th}}R_{\mathrm{th}}$ & 20分 \\
ヒーターで温める（6-4） & 温度 $T$ & $T_{\mathrm a}+PR_{\mathrm{th}}$ & $C_{\mathrm{th}}R_{\mathrm{th}}$ & \\
RC の充電（6-4） & 電圧 $V$ & $E$ & $RC$ & \qty{2}{s} \\
RL の電流（6-4・6-5） & 電流 $i$ & $E/R$ & $L/R$ & \\
落ちる物（6-5） & 速さ $v$ & $mg/b$ & $m/b$ & \\
2つの水そうの差（6-14） & 水位の差 & 0 & $\dfrac{1}{0.2}$ 分 & 5分 \\
\end{tblr}
\bktabcaption{1回微分の仲間。1つの答えの形で、全部が表せる。}
\end{bkcenter}

\kviphead{2回微分の仲間：ゆれの速さと減衰比}
2回微分の式は、相似の表で文字をおきかえるだけで、6-6・6-7 で作った式がそのまま使えます。
\begin{bkcenter}
\small
\begin{tblr}{width=\linewidth, colspec={Q[l,m,wd=30mm] X[c,m] X[c,m]}, row{1}={bg=bkTheme!12, font=\sffamily\small}, hline{1,2,Z}={0.6pt,bkGray}, hline{3-5}={0.3pt,bkGray}, rowsep=3pt, colsep=4pt}
 & ばね・質量・ダンパー & RLC 回路 \\
固有角振動数 $\omega_{\mathrm n}$ & $\sqrt{\dfrac{k}{m}}$ & $\dfrac{1}{\sqrt{LC}}$ \\
臨界の値 & $c_{\mathrm c}=2\sqrt{mk}$ & $R_{\mathrm c}=2\sqrt{\dfrac{L}{C}}$ \\
減衰比 $\zeta$ & $\dfrac{c}{2\sqrt{mk}}$ & $\dfrac{R}{2}\sqrt{\dfrac{C}{L}}$ \\
\end{tblr}
\bktabcaption{$k$ を $\frac1C$、$m$ を $L$、$c$ を $R$ にかえると、左の列から右の列ができる。}
\end{bkcenter}
たとえば、ばねの $\omega_{\mathrm n}=\sqrt{k/m}$ の $k$ に $\frac1C$、$m$ に $L$ を入れると、$\sqrt{\dfrac{1/C}{L}}=\dfrac{1}{\sqrt{LC}}$ になります。

\bksubsection{4}{同じ形なら、答えも同じ}

\label{bk:6-15:onaji}6-7 のばね（$m=\qty{0.5}{kg}$、$k=\qty{8}{N/m}$、$c=\qty{0.4}{N.s/m}$）と、RLC 回路（$L=\qty{10}{mH}$、$C=\qty{1}{\micro F}$、$R=\qty{20}{\ohm}$）をくらべます。
\begin{bkcenter}
\begin{bktable}{colspec={Q[l,m,wd=22mm] Q[c,m,wd=46mm] Q[c,m,wd=50mm]}, rowsep=2pt, row{1}={font=\sffamily\small}}
 & ばね & RLC 回路 \\
$\omega_{\mathrm n}$ & $\sqrt{8\div0.5}=4$\,rad/s & $\dfrac{1}{\sqrt{LC}}=\num{10000}$\,rad/s \\
$\zeta$ & $\dfrac{0.4}{2\sqrt{0.5\times8}}=0.1$ & $\dfrac{20}{2}\sqrt{\dfrac{0.000001}{0.01}}=0.1$ \\
ゆれの周期 & 約\qty{1.58}{s} & 約\qty{0.63}{ms} \\
\end{bktable}
\end{bkcenter}
減衰比 $\zeta$ が同じ0.1なので、ゆれ方の「形」はまったく同じです。ちがうのは、時間の物さしだけです。

\begin{bkfigure}
\begin{tikzpicture}
\begin{axis}[bkplot, axis x line=bottom, axis y line=left, width=60mm, height=40mm, xmin=0, xmax=8.4, ymin=-0.1, ymax=0.12, xtick={0,2,4,6,8}, ytick={-0.05,0,0.05,0.1}, xlabel={$t$（s）}, ylabel={$x$（m）}, yticklabel style={/pgf/number format/fixed}, title={ばね}, title style={font=\sffamily\small}]
  \addplot[bkthin, line width=0.5pt, domain=0:8.2] {0};
  \addplot[bkcurve, line width=1.2pt, domain=0:8.2, samples=300] {0.1*exp(-0.4*x)*(cos(deg(3.98*x))+0.1005*sin(deg(3.98*x)))};
\end{axis}
\end{tikzpicture}\hspace{6mm}%
\begin{tikzpicture}
\begin{axis}[bkplot, axis x line=bottom, axis y line=left, width=60mm, height=40mm, xmin=0, xmax=3.36, ymin=-10, ymax=12, xtick={0,0.8,1.6,2.4,3.2}, ytick={-5,0,5,10}, xlabel={$t$（ms）}, ylabel={電圧（V）}, title={RLC 回路}, title style={font=\sffamily\small}, xticklabel style={/pgf/number format/fixed}]
  \addplot[bkthin, line width=0.5pt, domain=0:3.3] {0};
  \addplot[bkcurve2, line width=1.2pt, domain=0:3.28, samples=300] {10*exp(-x)*(cos(deg(9.95*x))+0.1005*sin(deg(9.95*x)))};
\end{axis}
\end{tikzpicture}
\bkfigcaption{上：ばね（\qty{0.1}{m}引いてはなす）。下：RLC 回路のコンデンサーの電圧（\qty{10}{V}にためてつなぐ）。横軸の目もりは2500倍ちがうが、曲線の形はぴったり同じ。}
\end{bkfigure}

だから、1つの答えを作れば、全部に使えます。車のサスペンションを設計する人が「$\zeta$ をいくつにするか」と考えることと、電気の回路を設計する人が「抵抗をいくつにするか」と考えることは、数学としては同じ問題です。電気の回路は、部品をかえるのが簡単で、はかるのも速いので、むかしは機械の動きを電気の回路で調べました（コラム）。

\bksubsection{5}{エネルギーで見ると}

\label{bk:6-15:energy}運動方程式から、エネルギーのきまりを導いてみます。ばね $m\,\dd v/\dd t=-kx$ の両辺に、速度 $v=\dd x/\dd t$ をかけます。
\begin{bkhenkei}
mv\frac{\dd v}{\dd t}=-kx\frac{\dd x}{\dd t} & 両辺に $v$ をかけ、右辺は $v=\dfrac{\dd x}{\dd t}$ と書いた \\
\frac{\dd}{\dd t}\Bigl(\frac12mv^2\Bigr)=-\frac{\dd}{\dd t}\Bigl(\frac12kx^2\Bigr) & 合成関数の微分（5-6）を逆に使った \\
\frac{\dd}{\dd t}\Bigl(\frac12mv^2+\frac12kx^2\Bigr)=0 & 右辺を左へ移した \\
\end{bkhenkei}
2行めは、$\dfrac12mv^2$ を時間で微分すると $\dfrac12m\times2v\times\dfrac{\dd v}{\dd t}=mv\dfrac{\dd v}{\dd t}$ になること（$v^2$ を微分すると $2v$、そこに $v$ の微分をかける）を、右から左へ読んだものです。$\dfrac12kx^2$ も同じです。

微分して0なので、$\dfrac12mv^2+\dfrac12kx^2$（運動のエネルギー＋ばねのエネルギー）は一定です。5-12 で求めたばねの仕事 $\dfrac12kx^2$ が、ここで出てきました。運動の法則の中に、エネルギーが保たれるきまりがふくまれていたのです。

\begin{bkfigure}
\begin{tikzpicture}
\begin{axis}[bkplot, width=112mm, height=44mm, xmin=0, xmax=3.3, ymin=0, ymax=0.05, xtick={0.5,1,1.5,2,2.5,3}, ytick={0.02,0.04}, xlabel={$t$（s）}, ylabel={エネルギー（J）}, yticklabel style={/pgf/number format/fixed}, xticklabel style={/pgf/number format/fixed}]
  \addplot[bkcurve3, line width=1.3pt, domain=0:3.2] {0.04};
  \addplot[bkcurve, line width=1.1pt, domain=0:3.2, samples=200] {0.04*cos(deg(4*x))^2};
  \addplot[bkcurve2, line width=1.1pt, domain=0:3.2, samples=200] {0.04*sin(deg(4*x))^2};
  \node[font=\sffamily\footnotesize, anchor=south west, text=bkGreen!70!black] at (axis cs:2.35,0.0405) {合計（一定）};
\end{axis}
\end{tikzpicture}
\bkfigcaption{6-6 のばね（0.5\,kg、8\,N/m、\qty{0.1}{m}引いてはなす）。ばねのエネルギー $\frac12kx^2$（青）と運動のエネルギー $\frac12mv^2$（オレンジ）が入れかわり、合計はいつも\qty{0.04}{J}。}
\end{bkfigure}

ダンパーがあると、同じ計算で $\dfrac{\dd}{\dd t}\Bigl(\dfrac12mv^2+\dfrac12kx^2\Bigr)=-cv^2$ になります。$cv^2$ はいつも0以上なので、エネルギーはへり続けます。へった分は、ダンパーの中で熱になります（6-7⑥）。電気でも、$L\dfrac{\dd i}{\dd t}+Ri+\dfrac{q}{C}=0$ に $i$ をかけると、
\kvipline{$\dfrac{\dd}{\dd t}\Bigl(\dfrac12Li^2+\dfrac{q^2}{2C}\Bigr)=-Ri^2$}
となります。$\dfrac{q^2}{2C}=\dfrac12CV^2$ はコンデンサーのエネルギー（5-12）、$\dfrac12Li^2$ はコイルのエネルギー、$Ri^2$ は抵抗で1秒に熱になるエネルギーです。相似の表どおり、ばねとまったく同じ形です。

\bksubsection{6}{はじめの状態はいくつ要る？}

\label{bk:6-15:shoki}微分方程式だけでは、答えは1つに決まりません。はじめの状態が必要です（6-1）。いくつ必要かは、式の形で決まります。
\begin{bkcenter}
\begin{bktable}{width=\linewidth, colspec={X[l,m] Q[c,m,wd=26mm] X[l,m]}, rowsep=2pt, row{1}={font=\sffamily\small}}
式の形 & 必要な数 & 例 \\
1回微分の式 & 1つ & お湯のはじめの温度 $T(0)=80$\,℃ \\
2回微分の式 & 2つ & ばねのはじめの位置と速度 $x(0)$、$v(0)$ \\
1回微分の式が2つの連立 & 2つ & 2つの水そうのはじめの水位 $h_1(0)$、$h_2(0)$ \\
2回微分の式が2つの連立 & 4つ & 2つの台車の位置と速度（6-14） \\
\end{bktable}
\bktabcaption{はじめの状態の数は、積分するたびに出てくる定数の数と同じ（5-13）。}
\end{bkcenter}
時間でなく場所で変わる式では、はじめの状態のかわりに、両はしの状態（境界条件、6-3）で答えを決めることもあります。どちらの場合も、「ルール（微分方程式）」と「条件」の2つがそろって、はじめて答えがただ1つに決まります。この章の題の「微分方程式＋はじめの状態」は、このことです。

\bksubsection{7}{解き方の地図}

\label{bk:6-15:chizu}この巻で学んだ解き方を、1枚の地図にします。式を見たら、上から順に問いに答えていきます。
\begin{bkfigure}
% Ver2 指摘11・No.25：流れ図の組み方（FIX-ver2-common 2章）で組み直した。問いは左の1列、「いいえ」はまっすぐ下、「はい」はまっすぐ右。
%   すきま 4 mm 以上、矢印 0.9 pt・矢じり 2.6 mm、下への矢印は 10 mm、線は交わらない。問い②は、はい・いいえが1つに決まる文にした
\begin{tikzpicture}[x=1mm, y=1mm,
    q/.style={bk node, draw=bkOrange, fill=bkOrange!8, font=\sffamily\footnotesize, text width=42mm, minimum height=12mm, inner xsep=2mm, inner ysep=1.2mm},
    a/.style={bk node, font=\sffamily\footnotesize, text width=56mm, minimum height=12mm, inner xsep=2mm, inner ysep=1.2mm},
    fl/.style={-{Stealth[length=2.6mm, width=2.2mm]}, line width=0.9pt, bkInk!70, shorten <=0.5mm, shorten >=0.5mm},
    yes/.style={font=\sffamily\bfseries\footnotesize, text=bkInk, anchor=south, inner sep=0pt, yshift=1.2mm},
    no/.style={font=\sffamily\bfseries\footnotesize, text=bkInk, anchor=west, inner sep=0pt, xshift=1.5mm}]
  \node[q] (q1) at (0,0) {量が2つ以上からみあう？};
  \node[q] (q2) at (0,-24) {右辺は、時刻 $t$ と\\数だけでできている？};
  \node[q] (q3) at (0,-48) {1回微分で、係数が\\定数の線形な式？};
  \node[q] (q4) at (0,-72) {2回微分で、係数が\\定数の線形な式？};
  \node[a] (a1) at (70,0) {和と差・消去法（6-14）};
  \node[a] (a2) at (70,-24) {積分するだけ（6-3）};
  \node[a] (a3) at (70,-48) {目標へ近づく形（6-4・6-5）};
  \node[a] (a4) at (70,-72) {右辺0：特性方程式（6-6・6-12）\\右辺あり：ゆすられる（6-13）};
  \node[bk node, draw=bkPurple, fill=bkPurple!8, font=\sffamily\footnotesize, text width=46mm, inner xsep=2mm, inner ysep=1.5mm, anchor=north] (a5) at (0,-90)
    {線形でない・係数が変わる式\\表計算で解く（6-2）\\2乗の抵抗は $\tanh$（6-11）};
  \foreach \i in {1,2,3,4} {\draw[fl] (q\i.east) -- node[yes] {はい} (a\i.west);}
  \draw[fl] (q1.south) -- node[no] {いいえ} (q2.north);
  \draw[fl] (q2.south) -- node[no] {いいえ} (q3.north);
  \draw[fl] (q3.south) -- node[no] {いいえ} (q4.north);
  \draw[fl] (q4.south) -- node[no] {いいえ} (a5.north);
\end{tikzpicture}
\bkfigcaption{解き方の地図。どの道に進んでも、表計算（6-2）なら必ず数で解ける。式で解いたら、表計算の結果と重ねて確かめる。}
\end{bkfigure}
「線形」は、答えどうしをたしても答えになる式でした（5-18）。線形でも、$\dot y=-ty$ のように係数が時刻で変わる式は、6-4 や 6-12 の形にはなりません。$v^2$ や $xy$ のように、量どうしのかけ算が入ると線形でなくなり、式で解けることはまれです。そのときは、表計算の出番です。天気予報もロケットの軌道も、この道で計算されています（6-0）。

\Needspace{12\baselineskip}
\bksubsection{8}{この先へ：第6巻は乗りかえ駅}

\label{bk:6-15:sakie}第6巻の考えは、この先の巻と物理シリーズで、こう使われます。
\begin{bkcenter}
\begin{bktable}{width=\linewidth, colspec={Q[l,m,wd=40mm] X[l,m] Q[c,m,wd=20mm]}, rowsep=2pt, row{1}={font=\sffamily\small}}
第6巻の考え & この先で & 章 \\
連立の式を表にまとめる & 係数の表（行列）で、たくさんの量を一度に扱う。モードは固有ベクトル & 7-8・7-10・9-10 \\
場所でも変わる量 & 棒に広がる熱、弦を伝わる波。時間と場所の両方の微分が入る式（偏微分方程式） & 8-13 \\
$e^{st}$ と置く・$s$ の場所 & 微分を $s$ のかけ算に変える翻訳機（ラプラス変換）で、制御を設計する & 9-5・9-9 \\
相似の表 & 力学・熱・波動・電磁気の式を、同じ道具で解く & 物理シリーズ \\
\end{bktable}
\end{bkcenter}

\begin{taisetsu}
\def\tsl#1{\makebox[6.6\zw][l]{\sffamily\bfseries\small\color{\bkTT}#1}}
\tsl{物理の書き方}法則（微分方程式）＋はじめの状態 → 未来がただ1つに決まる\\[3pt]
\tsl{同じ形のわけ}ためるもの・慣性・失うもの・押すもの。きまりが比例で、合計でつなぐ\\[3pt]
\tsl{相似}力↔電圧、速度↔電流、$m\leftrightarrow L$、$c\leftrightarrow R$、$k\leftrightarrow\frac1C$、$x\leftrightarrow q$\\[3pt]
\tsl{エネルギー}運動方程式に $v$ をかけると出てくる。失う部品で熱になる\\[3pt]
\tsl{解き方}形を見て章を選ぶ。どれでも表計算なら解ける
\end{taisetsu}

\begin{kotoba}
\small
\bkword{相似}{ちがう分野のしくみが、同じ形の式で書けること。}
\bkword{キルヒホッフの電圧の法則}{回路を1周すると、電圧の上がり下がりの合計は0になるというきまり。}
\bkword{法則}{実験と観察で確かめられた、物理の出発点になるきまり（4-1）。}
\end{kotoba}

\par\addvspace{10pt}{\centering \bkkigocard{$\leftrightarrow$}{たいおうする}{相似の表で、2つの量が同じ役わりをすること。$m\leftrightarrow L$ のように書く}{corresponds to}{6-15}\par}\addvspace{12pt}

\begin{column}{アナログコンピュータ：微分方程式を電気で解く機械}
\noindent\begin{minipage}[c]{0.5\linewidth}
\raggedright
相似を使えば、機械の微分方程式を、同じ形の電気の回路を作って「はかる」ことで解けます。1931年、アメリカのブッシュは、歯車と回る円板で積分する「微分解析機」を作りました。1950〜60年代には、電子回路で積分するアナログコンピュータが、飛行機の設計などに使われました。この囲みの図のように、積分する部品を2つつなぎ、出てきた $x$ を $-\dfrac km$ 倍して入口にもどすと、回路の中で $\ddot x=-\dfrac km x$ がそのまま「動き」ます。今は、6-2 と同じ方法でコンピュータが計算しています。
\end{minipage}\hfill
\begin{minipage}[c]{0.46\linewidth}
\centering
\begin{tikzpicture}[x=1mm, y=1mm, b/.style={bk node, font=\sffamily\footnotesize, minimum width=12mm, minimum height=9mm, inner xsep=1mm}]
  \node[b] (i1) at (0,0) {積分};
  \node[b] (i2) at (28,0) {積分};
  \node[b, draw=bkOrange] (g) at (14,-18) {$\times\Bigl(-\dfrac km\Bigr)$};
  \draw[bk link] (i1) -- node[bk link label, above] {$\dot x$} (i2);
  \draw[bk link] (i2.east) -- ++(6,0) node[bk link label, above] {$x$} |- (g.east);
  \draw[bk link] (g.west) -| ([xshift=-6mm]i1.west) node[bk link label, above, pos=0.9] {$\ddot x$} -- (i1.west);
\end{tikzpicture}
\end{minipage}
\end{column}



%% ---------------- 例題 ----------------
\bksection{reidai}{例題でたしかめよう}

\begin{reidai}[法則から式を立てる：ばね・質量・ダンパー]
質量\qty{0.5}{kg}の台車が、ばね定数\qty{8}{N/m}のばねと、減衰係数\qty{0.4}{N.s/m}のダンパーにつながっています。台車を\qty{0.1}{m}引いて、静かにはなしました。運動の法則から微分方程式を立て、はじめの状態を書きましょう。また、どの章の方法で解けますか。\par
\kangaekata ②の手順。主役は位置 $x$ と速度 $v$。力は、ばね $-kx$ とダンパー $-cv$。2回微分の式なので、はじめの状態は2つ。\par
\shiki\kvips\hangindent=2.6\zw\hangafter=1 ばねの力：$-8x$、ダンパーの力：$-0.4v$（速度と反対向き）\par
\kvipa 運動の法則：$0.5\dfrac{\dd v}{\dd t}=-8x-0.4v$、\ $v=\dfrac{\dd x}{\dd t}$\par
\kvipa はじめの状態：引いた位置 $x(0)=0.1$、静かにはなしたので $v(0)=0$\par
\kvipa 2回微分で係数が定数の線形な式、右辺0 → 特性方程式（6-12）。表計算なら 6-7\par
\kotae\kvips\hangindent=2.6\zw\hangafter=1 $0.5\ddot x+0.4\dot x+8x=0$、$x(0)=0.1$、$v(0)=0$。6-12（6-7 の表計算）
\end{reidai}

\begin{jibun}
質量\qty{2}{kg}の台車が、ばね定数\qty{50}{N/m}のばねと減衰係数\qty{4}{N.s/m}のダンパーにつながっています。止まっている台車を、たたいて右向きに\qty{0.3}{m/s}の速さを与えました。微分方程式とはじめの状態を書きましょう。
\bkjans{$2\dfrac{\dd v}{\dd t}=-50x-4v$、$v=\dfrac{\dd x}{\dd t}$（$2\ddot x+4\dot x+50x=0$）。$x(0)=0$、$v(0)=0.3$。}
\end{jibun}

\begin{reidai}[法則から式を立てる：RC 回路]
電池（\qty{10}{V}）、抵抗（\qty{20}{k\ohm}）、空のコンデンサー（\qty{100}{\micro F}）をつなぎ、時刻0でスイッチを入れます。コンデンサーの電圧 $V$ の微分方程式と、はじめの状態、時定数、最後の電圧を求めましょう。\par
\kangaekata ②の「電池の電圧＝部品の電圧の合計」。抵抗の電圧は $Ri$、コンデンサーは $V$。電流 $i=C\dfrac{\dd V}{\dd t}$（$q=CV$ を微分）。\par
\shiki\kvips\hangindent=2.6\zw\hangafter=1 電圧の合計：$10=Ri+V$ より $i=\dfrac{10-V}{R}$\par
\kvipa 電流は電気のたまる速さ：$C\dfrac{\dd V}{\dd t}=\dfrac{10-V}{R}$\par
\kvipa $RC=\qty{20e3}{\ohm}\times\qty{100e-6}{F}=\qty{2}{s}$ なので $\dfrac{\dd V}{\dd t}=\dfrac{10-V}{2}$\par
\kvipa はじめは空なので $V(0)=0$。最後は変わらなくなる（$\dfrac{\dd V}{\dd t}=0$）ので $V=10$\par
\kotae\kvips\hangindent=2.6\zw\hangafter=1 $\dfrac{\dd V}{\dd t}=\dfrac{10-V}{2}$、$V(0)=0$、$\tau=\qty{2}{s}$、最後は\qty{10}{V}
\end{reidai}

\begin{jibun}
電池（\qty{5}{V}）、抵抗（\qty{10}{k\ohm}）、空のコンデンサー（\qty{47}{\micro F}）で、例題2と同じことを求めましょう。
\bkjans{$RC=\qty{10e3}{\ohm}\times\qty{47e-6}{F}=\qty{0.47}{s}$。$\dfrac{\dd V}{\dd t}=\dfrac{5-V}{0.47}$、$V(0)=0$、$\tau=\qty{0.47}{s}$、最後は\qty{5}{V}。}
\end{jibun}

\begin{reidai}[法則から式を立てる：冷めるお湯]
熱容量 $C_{\mathrm{th}}=\qty{1000}{J/K}$ のお湯（はじめ80\,℃）が、熱抵抗 $R_{\mathrm{th}}=\qty{1.2}{K/W}$ で20\,℃の部屋へ熱を逃がします。温度 $T$ の微分方程式（$t$ は秒）と時定数を求め、$t$ を分にした式に直しましょう。\par
\kangaekata ②の「たまる熱＝入る熱－出る熱」。入る熱は0、出る熱は $\dfrac{T-20}{R_{\mathrm{th}}}$。\par
\shiki\kvips\hangindent=2.6\zw\hangafter=1 $1000\dfrac{\dd T}{\dd t}=-\dfrac{T-20}{1.2}$ より $\dfrac{\dd T}{\dd t}=-\dfrac{T-20}{1200}$（1/s）\par
\kvipa 時定数 $\tau=C_{\mathrm{th}}R_{\mathrm{th}}=1000\times1.2=\qty{1200}{s}$\par
\kvipa 分に直す：1分は60秒なので、1分あたりの変わる速さは60倍\par
\kvipa $\dfrac{\dd T}{\dd t}=-\dfrac{60}{1200}(T-20)=-0.05(T-20)$（1/分）\par
\kotae\kvips\hangindent=2.6\zw\hangafter=1 $\dfrac{\dd T}{\dd t}=-\dfrac{T-20}{1200}$（秒）、$\tau=\qty{1200}{s}=20$分、$\dfrac{\dd T}{\dd t}=-0.05(T-20)$（分）、$T(0)=80$
\end{reidai}
6-1 から使ってきた冷めるお湯の「0.05」は、熱容量と熱抵抗から、こうして出てきます（$\qty{1200}{s}\div\qty{60}{s/min}=20$ 分）。

\begin{jibun}
熱容量\qty{1500}{J/K}のみそ汁が、熱抵抗\qty{0.5}{K/W}で25\,℃の部屋へ熱を逃がします。温度 $T$ の微分方程式（秒）と時定数を求め、分の式に直しましょう。
\bkjans{$1500\dfrac{\dd T}{\dd t}=-\dfrac{T-25}{0.5}$ より $\dfrac{\dd T}{\dd t}=-\dfrac{T-25}{750}$（秒）。$\tau=\qty{750}{s}=12.5$分。分の式は $\dfrac{\dd T}{\dd t}=-0.08(T-25)$。}
\end{jibun}

\begin{reidai}[相似の表で、答えを使い回す]
RLC 回路（$L=\qty{10}{mH}$、$C=\qty{1}{\micro F}$、$R=\qty{20}{\ohm}$）と、ばね・質量・ダンパー（$m=\qty{0.5}{kg}$、$k=\qty{8}{N/m}$、$c=\qty{0.4}{N.s/m}$）があります。それぞれの $\omega_{\mathrm n}$ と $\zeta$ を求め、ゆれ方の形が同じかを答えましょう。\par
\kangaekata ③の2回微分の表。ばねの式の $k$ を $\frac1C$、$m$ を $L$、$c$ を $R$ にかえる。\par
\shiki\kvips\hangindent=2.6\zw\hangafter=1 ばね：$\omega_{\mathrm n}=\sqrt{\dfrac{8}{0.5}}=4$、\ $\zeta=\dfrac{0.4}{2\sqrt{0.5\times8}}=\dfrac{0.4}{4}=0.1$\par
\kvipa 回路：$LC=0.01\times0.000001=10^{-8}$ なので $\omega_{\mathrm n}=\dfrac{1}{\sqrt{10^{-8}}}=10^{4}$\par
\kvipa \phantom{回路：}$\zeta=\dfrac{R}{2}\sqrt{\dfrac{C}{L}}=10\times\sqrt{\dfrac{0.000001}{0.01}}=10\times0.01=0.1$\par
\kvipa $\zeta$ が同じなので、ゆれ方の形は同じ。時間の物さしが $10^4\div4=2500$ 倍ちがう\par
\kotae\kvips\hangindent=2.6\zw\hangafter=1 ばね $\omega_{\mathrm n}=\qty{4}{rad/s}$、回路 $\omega_{\mathrm n}=\qty{10000}{rad/s}$、どちらも $\zeta=0.1$。形は同じ
\end{reidai}

\begin{jibun}
ばね・質量・ダンパー（$m=\qty{2}{kg}$、$k=\qty{200}{N/m}$、$c=\qty{4}{N.s/m}$）と、RLC 回路（$L=\qty{1}{mH}$、$C=\qty{10}{\micro F}$、$R=\qty{2}{\ohm}$）の $\omega_{\mathrm n}$ と $\zeta$ を求め、ゆれ方の形が同じかを答えましょう。
\bkjans{ばね：$\omega_{\mathrm n}=\sqrt{200\div2}=\qty{10}{rad/s}$、$\zeta=\dfrac{4}{2\sqrt{400}}=0.1$。回路：$LC=10^{-8}$ で $\omega_{\mathrm n}=\qty{10000}{rad/s}$、$\zeta=\dfrac{2}{2}\sqrt{\dfrac{0.00001}{0.001}}=0.1$。形は同じ（時間の物さしが1000倍ちがう）。}
\end{jibun}

\begin{reidai}[時定数をくらべる]
冷めるお湯（$\tau=20$ 分、はじめ80\,℃、部屋20\,℃）と、RC の充電（$\tau=\qty{2}{s}$、$E=\qty{10}{V}$、はじめ0\,V）があります。それぞれ「ほぼ終わる」（目標までの差が5\%になる）のは何時間後・何秒後ですか。そのときの温度と電圧も求めましょう（$e^{-3}=0.0498$）。\par
\kangaekata 目標との差は $e^{-t/\tau}$ 倍になる。$e^{-3}\approx0.05$ なので、$3\tau$ で差が5\%（6-4）。\par
\shiki\kvips\hangindent=2.6\zw\hangafter=1 お湯：$3\tau=3\times20=60$ 分。差は $60\times0.0498\approx2.99$\,℃ なので $T\approx20+2.99=23.0$\par
\kvipa RC：$3\tau=3\times2=\qty{6}{s}$。差は $10\times0.0498=0.498$\,V なので $V=10-0.498\approx9.50$\par
\kotae\kvips\hangindent=2.6\zw\hangafter=1 お湯は60分（1時間）後に約23.0\,℃、RC は6秒後に約\qty{9.50}{V}
\end{reidai}
同じ「$3\tau$ でほぼ終わり」が、お湯では1時間、電気では6秒です。式が同じ形なので、時定数さえわかれば、どちらも同じ考えで見積もれます。

\begin{jibun}
RL 回路（$L=\qty{0.5}{H}$、$R=\qty{10}{\ohm}$、$E=\qty{10}{V}$、はじめ電流0）と、冷めるみそ汁（$\tau=15$ 分、はじめ70\,℃、部屋25\,℃）が「ほぼ終わる」時刻と、そのときの電流・温度を求めましょう（$e^{-3}=0.0498$）。
\bkjans{RL：$\tau=\dfrac{L}{R}=\qty{0.05}{s}$、最後の電流 $\dfrac{E}{R}=\qty{1}{A}$。$3\tau=\qty{0.15}{s}$ で $1-0.0498\approx\qty{0.95}{A}$。みそ汁：$3\tau=45$ 分で $25+45\times0.0498\approx27.2$\,℃。}
\end{jibun}

\begin{reidai}[式の形から、解き方を選ぶ]
次の式は、⑦の地図のどの道で解けますか。章と、必要なはじめの状態の数を答えましょう。\par
\begin{tblr}{colspec={Q[l,m,wd=58mm] Q[l,m,wd=58mm]}, rowsep=1pt, colsep=0pt}
（1）$\dfrac{\dd y}{\dd t}=3t^2$ & （2）$\dfrac{\dd T}{\dd t}=-0.1(T-25)$ \\
（3）$2\ddot x+3\dot x+50x=0$ & （4）$\ddot x+9x=2\cos t$ \\
（5）$\dfrac{\dd v}{\dd t}=9.8-0.003v^2$ & \\
\end{tblr}\par
\kangaekata 地図の問いに上から答える。必要な数は⑥の表（1回微分は1つ、2回微分は2つ）。\par
\shiki\kvips\hangindent=2.6\zw\hangafter=1 （1）右辺は $t$ と数だけ → 6-3。1回微分なので1つ\par
\kvipa （2）1回微分・係数が定数の線形 → 6-4。1つ\par
\kvipa （3）2回微分・係数が定数の線形・右辺0 → 6-12。2つ\par
\kvipa （4）2回微分・係数が定数の線形・右辺あり → 6-13。2つ\par
\kvipa （5）$v^2$ があり線形でない → 6-2（6-11 の $\tanh$）。1つ\par
\kotae\kvips\hangindent=2.6\zw\hangafter=1 （1）6-3、1つ。（2）6-4、1つ。（3）6-12、2つ。（4）6-13、2つ。（5）6-2・6-11、1つ
\end{reidai}

\begin{jibun}
次の式は、どの章の方法で解けますか。必要なはじめの状態の数も答えましょう。\par
\begin{tblr}{colspec={Q[l,m,wd=50mm] Q[l,m,wd=66mm]}, rowsep=1pt, colsep=0pt}
（1）$\dfrac{\dd V}{\dd t}=5$ & （2）$0.5\ddot x=-8x$ \\
（3）$\dfrac{\dd N}{\dd t}=-0.2N$ & （4）$\dfrac{\dd h_1}{\dd t}=-0.1(h_1-h_2)$、$\dfrac{\dd h_2}{\dd t}=0.1(h_1-h_2)$ \\
\end{tblr}
\bkjans{（1）\nobreak 6-3、1つ。（2）\nobreak 6-6（6-12）、2つ。（3）\nobreak 6-4、1つ。（4）\nobreak 6-14、2つ。}
\end{jibun}

\begin{reidai}[エネルギーで速さ・電流を求める]
（1）6-6 のばね（$m=\qty{0.5}{kg}$、$k=\qty{8}{N/m}$）を\qty{0.1}{m}引いてはなしました。エネルギーの合計と、いちばん速いときの速さ、$x=\qty{0.05}{m}$ のときの速さを求めましょう。\par
（2）$C=\qty{1}{\micro F}$ のコンデンサーを\qty{10}{V}にため、$L=\qty{10}{mH}$ のコイルにつなぎました。いちばん大きい電流を求めましょう。\par
\kangaekata ⑤の「エネルギーの合計は一定」。速さがいちばん大きいのは $x=0$（ばねのエネルギーが0）のとき。電気も相似で同じ。\par
\shiki\kvips\hangindent=4.3\zw\hangafter=1 （1）合計：$\dfrac12\times8\times0.1^2=0.04$（はじめは $v=0$）\par
\kvipa \hspace*{1.7\zw}$x=0$：$\dfrac12\times0.5\times v^2=0.04$、$v^2=0.16$、$v=0.4$\par
\kvipa \hspace*{1.7\zw}$x=0.05$：$\dfrac12\times8\times0.05^2=0.01$ なので $\dfrac12\times0.5\times v^2=0.03$、$v=\sqrt{0.12}\approx0.346$\par
\kvipa （2）はじめ：$\dfrac12CV^2=\dfrac12\times0.000001\times10^2=0.00005$（J）\par
\kvipa \hspace*{1.7\zw}$q=0$ のとき全部コイルに：$\dfrac12\times0.01\times i^2=0.00005$、$i^2=0.01$、$i=0.1$\par
\kotae\kvips\hangindent=4.3\zw\hangafter=1 （1）\qty{0.04}{J}、\qty{0.4}{m/s}、約\qty{0.35}{m/s}。（2）\qty{0.1}{A}
\end{reidai}

\begin{jibun}
（1）$m=\qty{0.2}{kg}$、$k=\qty{20}{N/m}$ のばねを\qty{0.05}{m}引いてはなしました。エネルギーの合計と、いちばん速いときの速さを求めましょう。\newline （2）$C=\qty{4}{\micro F}$ を\qty{5}{V}にため、$L=\qty{25}{mH}$ につなぎました。いちばん大きい電流を求めましょう。
\bkjans{（1）\nobreak $\dfrac12\times20\times0.05^2=\qty{0.025}{J}$。$\dfrac12\times0.2\times v^2=0.025$ より $v=\qty{0.5}{m/s}$。（2）\nobreak $\dfrac12\times0.000004\times25=0.00005$\,J。$\dfrac12\times0.025\times i^2=0.00005$ より $i^2=0.004$、$i\approx\qty{0.063}{A}$。}
\end{jibun}

%% ---------------- まちがえやすいところ ----------------
\bksection{tsumazuki}{まちがえやすいところをチェック}

\noindent\begin{minipage}[t]{0.485\linewidth}\vspace{0pt}
\begin{tsumazuki}{① 相似の組を逆にする}
\bkng{力 $\leftrightarrow$ 電流}
\bkok{力 $\leftrightarrow$ 電圧、速度 $\leftrightarrow$ 電流}
ばねの式とRLCの式をならべて、同じ場所にある量を組にします。$F$ と $E$ が右辺にならびます。
\end{tsumazuki}
\end{minipage}\hfill
\begin{minipage}[t]{0.485\linewidth}\vspace{0pt}
\begin{tsumazuki}{② ばね定数とコンデンサー}
\bkng{$k\leftrightarrow C$}
\bkok{$k\leftrightarrow\dfrac1C$（$\dfrac1k\leftrightarrow C$）}
ばねの力 $kx$ とコンデンサーの電圧 $\dfrac1C q$ をくらべます。
\end{tsumazuki}
\end{minipage}

\noindent\begin{minipage}[t]{0.485\linewidth}\vspace{0pt}
\begin{tsumazuki}{③ 熱もゆれると思う}
\bkng{お湯の温度も、部屋の温度のまわりでゆれる}
\bkok{熱には慣性がないので、ゆれずに近づく}
相似の表の「慣性」の行が、熱では「ない」からです。
\end{tsumazuki}
\end{minipage}\hfill
\begin{minipage}[t]{0.485\linewidth}\vspace{0pt}
\begin{tsumazuki}{④ はじめの状態が足りない}
\bkng{ばね：$x(0)=0.1$ だけ}
\bkok{$x(0)=0.1$、$v(0)=0$}
2回微分の式には、はじめの位置と速度の2つが要ります。「静かにはなす」は $v(0)=0$ のことです。
\end{tsumazuki}
\end{minipage}

%% ---------------- 物理・工学ではこう使う ----------------
\begin{butsuri}
\kvipbox{パソコンの CPU の熱設計}
\noindent\begin{minipage}[c]{0.58\linewidth}
\raggedright
パソコンの CPU（計算をする半導体のチップ）は、動くと熱を出します。設計する人は、熱の列を電気の回路のように考えます。発熱\qty{100}{W}、チップから空気までの熱抵抗が\qty{0.3}{K/W}なら、最後の温度は空気より $100\times0.3=30$\,K 高くなります（6-4 の $T_{\mathrm a}+PR_{\mathrm{th}}$。数はだいたいの例）。放熱器の熱容量が大きいほど時定数が長く、温度はゆっくり上がります。
\end{minipage}\hfill
\begin{minipage}[c]{0.38\linewidth}
\centering
\begin{tikzpicture}[x=1mm, y=1mm]
  \fill[bkGray!40] (4,0) rectangle (40,2);
  \fill[bkInk!70] (14,2) rectangle (30,5);
  \foreach \x in {8,13,...,38} {\fill[bkBlue!35] (\x,8) rectangle ({\x+2.4},24);}
  \fill[bkBlue!35] (6,5) rectangle (38,8);
  \draw[bkBlue!70!black, line width=0.6pt] (6,5) rectangle (38,8);
  \draw[bkOrange, line width=1pt, -{Stealth[length=3pt]}] (22,10) -- (22,28);
  \node[font=\sffamily\footnotesize, anchor=west] at (31,3.5) {CPU};
  \node[font=\sffamily\footnotesize, anchor=south] at (22,28.5) {熱 $P$};
  \node[font=\sffamily\footnotesize, anchor=north] at (22,-1) {放熱器（ヒートシンク）};
\end{tikzpicture}
\end{minipage}
\end{butsuri}

\begin{butsuri}
\kvipbox{地震計：ばね・質量・ダンパーで地面のゆれをはかる}
\noindent\begin{minipage}[c]{0.58\linewidth}
\raggedright
むかしの地震計は、ばねでつるしたおもりと、ダンパーでできていました。地面が速くゆれると、おもりは慣性でほとんど動かず、地面だけが動きます。地面とおもりのずれを記録すれば、地面のゆれがわかります。おもりが自分でゆれ続けないように、ダンパーで $\zeta$ を0.7くらいにします（6-12⑦）。
\end{minipage}\hfill
\begin{minipage}[c]{0.38\linewidth}
\centering
\begin{tikzpicture}[x=1mm, y=1mm]
  \fill[bkInk!40] (0,0) rectangle (40,2);
  \draw[bkInk!80, line width=1pt] (4,2) -- (4,30) -- (36,30) -- (36,2);
  \draw[bkInk, line width=0.8pt, decorate, decoration={zigzag, amplitude=1.3mm, segment length=2.2mm, pre length=1.5mm, post length=1.5mm}] (20,30) -- (20,18);
  \filldraw[fill=bkBlue!25, draw=bkBlue!80!black, line width=0.8pt, rounded corners=0.8mm] (14,10) rectangle (26,18);
  \node[font=\sffamily\small] at (20,14) {$m$};
  \draw[bkInk, line width=0.8pt] (20,10) -- (20,7) (17,7) -- (23,7) (17,3) -- (17,7.5) (23,3) -- (23,7.5) (20,2) -- (20,4.6);
  \draw[bkInk, line width=1.4pt] (18.4,4.6) -- (21.6,4.6);
  \draw[bkRed, line width=1pt, {Stealth[length=3pt]}-{Stealth[length=3pt]}] (2,-3) -- (38,-3);
  \node[font=\sffamily\footnotesize, anchor=north] at (20,-4) {地面のゆれ};
\end{tikzpicture}
\end{minipage}
\end{butsuri}

%% ---------------- 章末問題 ----------------
\bkendproblems
\bkproblemgroup{基本}{2}

\bkproblem 質量\qty{1}{kg}の台車が、ばね定数\qty{25}{N/m}のばねと減衰係数\qty{2}{N.s/m}のダンパーにつながっています。\qty{0.2}{m}引いて静かにはなしました。微分方程式とはじめの状態を書きましょう\mbox{（\bkreidairef{1}）}。
\bkpans{$1\times\dfrac{\dd v}{\dd t}=-25x-2v$、$v=\dfrac{\dd x}{\dd t}$（$\ddot x+2\dot x+25x=0$）。$x(0)=0.2$、$v(0)=0$。}

\bkproblem 電池（\qty{9}{V}）、抵抗（\qty{100}{k\ohm}）、空のコンデンサー（\qty{10}{\micro F}）をつなぎます。コンデンサーの電圧の微分方程式、はじめの状態、時定数、最後の電圧を求めましょう\mbox{（\bkreidairef{2}）}。
\bkpans{$RC=\qty{100e3}{\ohm}\times\qty{10e-6}{F}=\qty{1}{s}$。$\dfrac{\dd V}{\dd t}=9-V$、$V(0)=0$、$\tau=\qty{1}{s}$、最後は\qty{9}{V}。}

\bkproblem 熱容量\qty{900}{J/K}の飲み物が、熱抵抗\qty{2}{K/W}で22\,℃の部屋へ熱を逃がします。温度の微分方程式（秒）と時定数を求め、分の式に直しましょう\mbox{（\bkreidairef{3}）}。
\bkpans{$900\dfrac{\dd T}{\dd t}=-\dfrac{T-22}{2}$ より $\dfrac{\dd T}{\dd t}=-\dfrac{T-22}{1800}$（秒）。$\tau=\qty{1800}{s}=30$分。分の式は $\dfrac{\dd T}{\dd t}=-\dfrac{1}{30}(T-22)\approx-0.033(T-22)$。}

\bkproblem RLC 回路（$L=\qty{0.1}{H}$、$C=\qty{10}{\micro F}$、$R=\qty{40}{\ohm}$）の $\omega_{\mathrm n}$ と $\zeta$ を求めましょう。また、相似の表で、この回路と同じ形のばね・質量・ダンパーでは、$L$、$R$、$\frac1C$ が何に対応しますか\mbox{（\bkreidairef{4}）}。
\bkpans{$LC=10^{-6}$ で $\omega_{\mathrm n}=\qty{1000}{rad/s}$。$\zeta=\dfrac{40}{2}\sqrt{\dfrac{0.00001}{0.1}}=20\times0.01=0.2$。$L\leftrightarrow m$（質量）、$R\leftrightarrow c$（ダンパー）、$\frac1C\leftrightarrow k$（ばね定数）。}

\bkproblem RC の充電（$R=\qty{5}{k\ohm}$、$C=\qty{100}{\micro F}$、$E=\qty{6}{V}$、はじめ0\,V）が「ほぼ終わる」（差が5\%）時刻と、そのときの電圧を求めましょう\mbox{（\bkreidairef{5}）}（$e^{-3}=0.0498$）。
\bkpans{$\tau=\qty{5e3}{\ohm}\times\qty{100e-6}{F}=\qty{0.5}{s}$。$3\tau=\qty{1.5}{s}$。$V=6-6\times0.0498\approx\qty{5.70}{V}$。}

\bkproblem 次の式は、どの章の方法で解けますか。必要なはじめの状態の数も答えましょう\mbox{（\bkreidairef{6}）}。\newline （1）$\dfrac{\dd x}{\dd t}=2t+1$\newline （2）$\ddot x+4\dot x+4x=0$\newline （3）$\dfrac{\dd v}{\dd t}=9.8-0.5v$\newline （4）$\ddot x+4x=3\cos t$
\bkpans{（1）\nobreak 6-3、1つ。（2）\nobreak 6-12、2つ。（3）\nobreak 6-4・6-5、1つ。（4）\nobreak 6-13、2つ。}

\bkproblemgroup{チャレンジ}{3}

\bkproblem （1）$m=\qty{1}{kg}$、$k=\qty{100}{N/m}$ のばねを\qty{0.03}{m}引いてはなしました。エネルギーの合計と、いちばん速いときの速さを求めましょう。\newline （2）$C=\qty{2}{\micro F}$ を\qty{10}{V}にため、$L=\qty{50}{mH}$ につなぎました。いちばん大きい電流を求めましょう\mbox{（\bkreidairef{7}）}。
\bkpans{（1）\nobreak $\dfrac12\times100\times0.03^2=\qty{0.045}{J}$。$\dfrac12\times1\times v^2=0.045$ より $v=\qty{0.3}{m/s}$。（2）\nobreak $\dfrac12\times0.000002\times100=0.0001$\,J。$\dfrac12\times0.05\times i^2=0.0001$ より $i^2=0.004$、$i\approx\qty{0.063}{A}$。}

\bkproblem 熱容量\qty{1050}{J/K}のお湯（250\,mL）が、熱抵抗\qty{1.1}{K/W}で部屋へ熱を逃がします。時定数は何分ですか。また、「ほぼ冷めきる」（差が5\%になる）のは約何分後ですか\mbox{（\bkreidairef{3}・\bkreidairef{5}）}。
\bkpans{$\tau=1050\times1.1=\qty{1155}{s}$。$\qty{1155}{s}\div\qty{60}{s/min}\approx19$分。$3\tau\approx58$分後。}

\bkproblem ばね・質量・ダンパー（$m=\qty{1}{kg}$、$k=\qty{400}{N/m}$、$c=\qty{8}{N.s/m}$）と、RLC 回路（$L=\qty{2}{mH}$、$C=\qty{5}{\micro F}$、$R=\qty{8}{\ohm}$）の $\omega_{\mathrm n}$ と $\zeta$ を求め、ゆれ方の形が同じかを答えましょう\mbox{（\bkreidairef{4}）}。
\bkpans{ばね：$\omega_{\mathrm n}=\sqrt{400}=\qty{20}{rad/s}$、$\zeta=\dfrac{8}{2\times20}=0.2$。回路：$LC=10^{-8}$ で $\omega_{\mathrm n}=\qty{10000}{rad/s}$、$\zeta=\dfrac{8}{2}\sqrt{\dfrac{0.000005}{0.002}}=4\times0.05=0.2$。形は同じ（時間の物さしが500倍ちがう）。}

\bkproblem RC の充電（$R=\qty{1}{k\ohm}$、$C=\qty{1000}{\micro F}$、$E=\qty{12}{V}$、はじめ0\,V）と、冷めるお湯（$\tau=25$ 分、はじめ80\,℃、部屋20\,℃）が「ほぼ終わる」時刻と、そのときの電圧・温度を求めましょう\mbox{（\bkreidairef{5}）}（$e^{-3}=0.0498$）。
\bkpans{RC：$\tau=\qty{1e3}{\ohm}\times\qty{1000e-6}{F}=\qty{1}{s}$、$3\tau=\qty{3}{s}$ で $12-12\times0.0498\approx\qty{11.4}{V}$。お湯：$3\tau=75$ 分で $20+60\times0.0498\approx23.0$\,℃。}

\begin{furikaeri}
\begin{checklist}
\item 法則から、微分方程式とはじめの状態を書ける。\bkfurref{bk:6-15:buhin}
\item 相似の表で、ばね・電気・熱の量を対応させられる。\bkfurref{bk:6-15:soji}
\item 式の形から、解き方（章）を選べる。\bkfurref{bk:6-15:chizu}
\end{checklist}
\end{furikaeri}

\bkfukidashi{delta}{第6巻、おつかれさま！　ルールとはじめの状態があれば、未来は計算できる。第7巻では、たくさんの量を表のまま計算しよう。}
